% Generated draft: requires mathematical and source review.
\documentclass[11pt]{article}
\usepackage[margin=25mm]{geometry}
\usepackage{fontspec}
\setmainfont{Latin Modern Roman}
\usepackage{amsmath,amssymb,mathtools,unicode-math}
\setmathfont{Latin Modern Math}
\usepackage{enumitem,needspace,xurl,hyperref}
\setlength{\parindent}{0pt}
\setlength{\parskip}{5pt}
\setlength{\emergencystretch}{3em}
\newcommand{\N}{\mathbb{N}}
\newcommand{\Z}{\mathbb{Z}}
\newcommand{\Q}{\mathbb{Q}}
\newcommand{\R}{\mathbb{R}}
\newcommand{\C}{\mathbb{C}}
\newcommand{\F}{\mathbb{F}}
\newcommand{\A}{\mathbb{A}}
\newcommand{\PP}{\mathbb{P}}
\newcommand{\E}{\mathbb{E}}
\newcommand{\eps}{\varepsilon}
\newcommand{\dd}{\,\mathrm{d}}
\begin{document}
\section*{UnsolvedMath ID: 800460}
\textbf{Source problem code:} catalogue entry
\subsection{Definitions and mathematical statement}
The supplied record leaves several conventions unstated and writes an undefined target \(M\). The following is a conditional formulation using the standard interpretations described here; these interpretations require verification against the source.

Let \(n\geq 1\), and let \((M_i^n,g_i,p_i)\) be pointed \(n\)-dimensional Riemannian manifolds. Interpret the displayed convergence as pointed Gromov--Hausdorff convergence
\[
(M_i^n,d_{g_i},p_i)\longrightarrow (X,d,p).
\]
Assume that there are constants \(\lambda\geq 0\) and \(v>0\), independent of \(i\), such that
\[
\operatorname{Ric}_{g_i}\geq -\lambda g_i,
\qquad
\operatorname{Vol}_{g_i}\bigl(B_1(p_i)\bigr)>v
\quad\text{for every }i.
\]
Here \(B_1(p_i)\) is the open metric unit ball. Completeness of the manifolds and the precise convergence convention are not explicitly stated in the supplied excerpt.

A tangent cone at \(x\in X\) means a pointed Gromov--Hausdorff limit of \((X,r_j^{-1}d,x)\) along a sequence \(r_j\downarrow 0\). Under the standard metric-cone convention, a pointed space is \(\ell\)-symmetric if it is pointed-isometric to
\[
\R^\ell\times C(Z)
\]
with its product metric and basepoint \((0,o)\), where \(C(Z)\) is a metric cone with vertex \(o\), allowing the residual cone to be a point. For a nonnegative integer \(k\), interpret the stratum as
\[
\mathcal S^k(X)
=
\{x\in X: \text{no tangent cone at }x\text{ is }(k+1)\text{-symmetric}\}.
\]
These definitions are not supplied in the excerpt, and its intended range of \(k\) is unspecified.

Subject to these conventions, and conditional on replacing the undefined target \(M\) in the source by \(X\), the conjecture asserts that there is an at most countable collection of bi-H\"older maps
\[
\varphi_j:B_1(0^k)\subset\R^k\longrightarrow X
\]
such that, on setting
\[
\widehat{\mathcal S}^{\,k}(X)
:=\bigcup_j\varphi_j\bigl(B_1(0^k)\bigr),
\]
one has
\begin{enumerate}
\item \(\mathcal S^k(X)\subseteq\widehat{\mathcal S}^{\,k}(X)\);
\item for every \(x\in\widehat{\mathcal S}^{\,k}(X)\), every tangent cone at \(x\) is at least \(k\)-symmetric, meaning that it admits a splitting with an \(\R^k\) factor as above.
\end{enumerate}
Here bi-H\"older is interpreted as a homeomorphism onto its image for which both the map and its inverse are H\"older continuous, using Euclidean distance on the domain and the restricted metric \(d\) on the image. The source specifies no exponents, constants, or uniformity requirements. Each image is consequently an embedded topological \(k\)-manifold without boundary; the source's phrase ``\(k\)-dimensional submanifolds'' does not specify additional smooth or locally flat structure. No such additional structure is imposed in this conditional formulation.
\subsection{Short English statement}
The supplied conjecture asks whether each stratum \(\mathcal S^k(X)\) of a noncollapsed Ricci limit space can be covered by countably many bi-H\"older images of open \(k\)-balls, with every tangent cone at every point of the covering admitting an \(\R^k\) factor. The target \(M\), the intended index range, and several background conventions remain unresolved; the displayed formulation conditionally uses target \(X\).
\subsection{Sources}
\begin{itemize}
\item Local UnsolvedMath registry: permanent ID \(800460\), exact source problem code \texttt{catalogue entry}, title ``Spaces with Lower Ricci Curvature --- Conjecture 4.'' The supplied record attributes the conjecture to Naber, with proposed year \(2020\), and labels it \texttt{open}; this is a supplied status, not an independently verified mathematical or literature assessment.
\item Naber, \emph{Conjectures and Open Questions on Spaces with Lower Ricci Curvature} (2020), Conjecture 4. Supplied source URL: \url{https://arxiv.org/abs/2010.10031}. The background records access on \texttt{2026-10-16} and extraction method \texttt{arxiv\_latex}. This access information is supplied provenance and does not indicate independent source consultation for this draft.
\item Supplied upstream ID: \(800460\); upstream revision: \texttt{d371ee917ac9647fbb06cd6fc529deec6f3dda82}. The local record was created and updated on \texttt{2026-10-16T14:00:00Z}.
\end{itemize}
\end{document}
